%%%%%%%%%%%%%%%%%%%%%%%%%%%%%%%%%%%%%%%%%%%%%%%%%%%%%%%%%%%%%%%%%%
%% Toto je inkludovaný LaTeXový soubor `cp_other.tex'.          %%
%% ------------------------------------------------------------ %%
%% Verze 2022-05-23                                             %%
%% Vyvíjí & udržuje <stanislav.hledik@physics.slu.cz>           %%
%%%%%%%%%%%%%%%%%%%%%%%%%%%%%%%%%%%%%%%%%%%%%%%%%%%%%%%%%%%%%%%%%%

\chapter{Některá drobná vylepšení}\label{improv}%%%%%%%%%%%%%%%%%%%%%%%%%%%

Pro vysázení motta kapitoly (sekce) lze použít prostředí \pgm{motto} s~jedním
povinným argumentem (autorem motta) a jedním volitelným argumentem
(specifikací jazyka motta). Příklad použití je v~začátku Úvodu na
stránce~\pageref{intro}.

Chcete-li vysadit poznámku, jež není součástí textu (např. pro vedoucího
práce, upomínku na dodatečné zjištění přesných bibliografických dat apod.),
použijte příkaz \verb!\note! (viz též stránka~\pageref{opdraft}). Argument
příkazu \verb!\note! se vysadí jako na ukázce na stránce~\pageref{opdraft},
ale jen při použití volby \opt{draft}. Jakmile ji vypneme, všechny marginálie
vysázené pomocí \verb!\note! zmizí a příkaz ani jeho argumenty není třeba ze
zdrojového textu odstraňovat. Příkaz \verb!\note! má hvězdičkovou variantu
\verb!\note*!, která vysází marginálii stejným způsobem jako \verb!\note!
s~tím rozdílem, že vypnutí volby \opt{draft} marginálii ponechá. Tímto
způsobem je vysázena marginálie na stránce~\pageref{opdraft}, která je
součástí finální formy manuálu.
% Tímto způsobem jsou vysázeny marginálie na stránkách~\pageref{opdraft}
% a~\pageref{mathant}, které jsou součástí finální formy manuálu.

\endinput

\section{Fragmentace zdrojového kódu}\label{fragment}%%%%%%%%%%%%%%%%%%%%%%

\section{Základy vytváření tabulek}\label{tables}%%%%%%%%%%%%%%%%%%%%%%%%%%

\section{Tvorba rejstříku}\label{treatindex}%%%%%%%%%%%%%%%%%%%%%%%%%%%%%%%

\section{Tipy pro některá IDE pro Windows}\label{winide}%%%%%%%%%%%%%%%%%%%

Někdy se může stát, že v~hypertextové \designation{pdf} verzi (viz
stránku~\pageref{pkghyperref}) vyžaduje v~některých místech odlišný zdrojový
text než tisknutelná verze; typicky se tato situace může vyskytnout při
odkazování na online zdroj (v~hypertextové verzi) nebo na knižní či
časopisecký zdroj (v~tištěné verzi). To lze jednoduše vyřešit pomocí
primitivního registru formátu \pkg{pdflatex} \verb!\pdfoutput!, jehož hodnota
je \texttt{0} pro \designation{dvi} výstup a \texttt{1} pro \designation{pdf}
výstup:

{\small
\begin{verbatim}
\ifcase\pdfoutput % (=0)
% sem přijde kód pro DVI nebo tisknutelné PDF, např.
(viz~\cite{Wil-etal:LIVE:PlottingProgram:Gnuplot})
\else % (>=1)
% sem přijde alternativní kód pro hypertextové PDF, např.
(viz \url{http://www.gnuplot.info/})
\fi
\end{verbatim}
\par}

%%
%% End of file `cp_other.tex'.
